\subsection{Integer interpolation and the quasi-separated positive-power theorem}

The regular-surface test admits a direct argument which extends further.
The key input is that powers of the given section lie in one affine
identity chart for infinitely many integers. This controls rational
coordinate denominators without constructing a relative Chevalley
quotient or assuming that a schematic closure is flat. The conclusion
first concerns the positive-power assertion in \S50.2, folios
652 verso--654, C4 scans 10--13. We then resolve the smooth-model
question of \S52.3, scan 213, while distinguishing its preceding
requirements from the supplementary axioms introduced afterwards.

\begin{statement}{D.17}{Proved}{Rational interpolation with an integral denominator}
Let $R$ be a Noetherian domain containing $\Q$, with fraction field
$K$. If $F\in K(T)$ is defined and belongs to $R$ at infinitely
many integers $T=n$, then
\[
       F=A(T)/B(T),\qquad A,B\in R[T],\qquad B(n_0)=1
\]
for some integer $n_0$, after multiplying numerator and denominator
by a sign if necessary.

If $R$ is local and $M$ is a finite torsion-free $R$-module, the
same argument for $F\in M_K\otimes_K K(T)$ with $F(n)\in M$
at infinitely many integers supplies
$F=N(T)/D(T)$, where $N\in M[T]$, $D\in R[T]$, and the
reduction of $D$ modulo the maximal ideal is nonzero. In particular
$D(n)$ is a unit at all but finitely many integers.
\end{statement}
\begin{proof}
We first prove the scalar assertion for a Krull domain $R$ containing
$\Q$; Noetherianity is not needed in this first step. Choose
$a,b\in R[T]$ representing $F$, with $a,b$ coprime over $K$.
Clearing a Bezout identity over $K[T]$ gives
\[
              u(T)a(T)+v(T)b(T)=r\ne0,
              \qquad u,v\in R[T],\ r\in R.
\]
On the infinite set of good integers, $b(n)\ne0$ and
$b(n)$ divides $a(n)$ in $R$, hence divides $r$.

Use the discrete valuations defining the Krull domain, with their
finite-character property. Only finitely many of them, say
$v_1,\ldots,v_h$, are positive on $r$. For each $v_i$, let $b_i$
be the minimum of the valuations of the nonzero coefficients of
$b$. Dividing by a uniformizer to this minimum and reducing gives
a nonzero polynomial over the residue field. Its values at
distinct integers are distinct inputs, since every nonzero
integer is a unit in $R$ and in these valuation rings. Thus
\[
                       v_i(b(n))=b_i
\]
except at finitely many integers. At every other defining
valuation, divisibility $b(n)\mid r$ forces $v(b(n))=0$.
After finitely many exclusions from the good integer set, all
$b(n)$ therefore have the same valuations. Fix one such $n_0$.
The Krull intersection description shows that
$b(n)/b(n_0)$ and its inverse lie in $R$ for every remaining
good $n$. Consequently both polynomials
\[
                        a/b(n_0),\qquad b/b(n_0)
\]
have values in $R$ at infinitely many integers. Vandermonde
interpolation at sufficiently many of them puts all their
coefficients in $R$: every nonzero integer difference is a unit.
The second polynomial has value $1$ at $n_0$.
The finite-character and intersection properties used here are
the definition and Lemma 1.1 of \citet[p.~998]{DS_Nagata1996}.

For a general Noetherian domain $R$, its integral closure
$\overline R\subset K$ is a Krull domain by the Mori--Nagata
theorem, in the precise form stated in
\citet[Lemma 1.2, p.~998]{DS_Nagata1996}. The preceding argument
gives $F=a'/b'$ with $a',b'\in\overline R[T]$ and
$b'(n_0)=1$. We now descend the denominator; finiteness of
$\overline R/R$ is not assumed.

Put $c=(b'-1)/(T-n_0)\in\overline R[T]$. This element is
integral over $R[T]$, since it has only finitely many coefficients
integral over $R$. Choose a monic relation
\[
 c^d+a_{d-1}(T)c^{d-1}+\cdots+a_0(T)=0,
                       \qquad a_i\in R[T].
\]
The polynomial
\[
 Q(U)=(U-1)^d+\sum_{i=0}^{d-1}
                   a_i(T)(T-n_0)^{d-i}(U-1)^i
\]
satisfies $Q(b')=0$. Write $Q(U)=UQ_1(U)+D(T)$. Then
\[
 D(n_0)=(-1)^d,\qquad D/b'=-Q_1(b')\in\overline R[T].
\]
Hence $F=N/D$ for $N=-a'Q_1(b')\in\overline R[T]$.
At all good integers outside a finite set,
$N(n)=F(n)D(n)\in R$. A second Vandermonde interpolation
puts $N$ in $R[T]$. Multiplying by $(-1)^d$ normalizes the
denominator at $n_0$. This proves the scalar assertion even
when normalization is not finite.

For the module assertion, choose an embedding $M\hookrightarrow R^q$
by choosing a basis of $M_K$ and clearing the coordinates of a
finite set of generators. Apply the scalar assertion to all
coordinates of $F$. Each resulting denominator has a unit
value somewhere, and thus a nonzero reduction over the residue
field of the local ring. Their product $D$ also has nonzero
reduction. Now $DF$ is a polynomial with coefficients in $R^q$,
and its values at infinitely many good integers belong to $M$.
Vandermonde interpolation expresses every coefficient as an
$R$-linear combination of those values, so every coefficient
belongs to $M$. This proves the module assertion.
\end{proof}

\begin{statement}{D.18}{Proved}{All Noetherian bases in characteristic zero}
Let $S$ be a Noetherian $\Q$-scheme, $G/S$ a separated group scheme
of finite presentation, and $g\in G(S)$ a fibrewise quasi-unipotent
section. There exist $m>0$ and a homomorphism
$u:\mathbb G_{a,S}\to G$ with $u(1)=g^m$.
For any fixed value at $1$, the homomorphism is unique.
The base may be reducible or nonreduced, and $G$ need not be
flat, smooth, or affine.
\end{statement}
\begin{proof}
We first record the denominator argument for a Noetherian local
$\Q$-algebra $A$ whose zero ideal is primary. Its nilradical
$\mathfrak p$ is prime, and every element outside $\mathfrak p$
is a nonzerodivisor. Thus the localization
\[
 B=(A\setminus\mathfrak p)^{-1}A
\]
contains $A$ injectively and is an Artinian local ring. Write
$J=\operatorname{rad}(B)$ and $R=A/\mathfrak p$, whose fraction
field is $K=B/J$. The Artinian ring $B$ has a coefficient field
$K$ by the equicharacteristic-zero case of the Cohen structure
theorem \citep[Theorem 10.160.8, Tag 032A]{DS_StacksCohen}.
It is finite-dimensional over that field.

For $j\ge1$, set $I_j=A\cap J^j\subset B$ and
$M_j=I_j/I_{j+1}$. These are finite modules, and
\[
 M_j\hookrightarrow J^j/J^{j+1},\qquad
 M_j\otimes_R K=J^j/J^{j+1}.
\]
The action factors through $R$ because
$\mathfrak p I_j\subset I_{j+1}$. The displayed inclusion
therefore makes $M_j$ a finite torsion-free $R$-module. The
filtration is saturated by its definition; one must not replace
it without justification by the potentially nonsaturated powers
of $\mathfrak p$ in $A$.

Let $B(T)=B\otimes_K K(T)$, and suppose $F\in B(T)$ has
values in $A$ at infinitely many integers. Reduce modulo $J$
and apply D.17 over the local domain $R$. Lift its numerator
and denominator to $A[T]$. The denominator is a unit at all
but finitely many integers. Subtracting the lifted fraction
leaves an element of $J B(T)$ whose good integer values belong
to $I_1$. Reduce the error modulo $J^2$ and apply the module
assertion of D.17 to $M_1$. Lift that numerator to $I_1[T]$
and its denominator to $A[T]$, and subtract again. Iteration
through the finite nilpotent filtration gives
\[
              F=N(T)/D(T),\qquad N,D\in A[T],
\]
where $D$ has nonzero reduction in the residue field of $A$.
Only finitely many integers were discarded at each step.
Products of the finitely many denominators still have nonzero
reduction. This proves the required denominator lemma for
primary rings, including their nilpotents.

Now let $G/A$ and $g$ be as in the statement, with $A$ primary
and local. After a positive power, assume that $g$ is unipotent
both over $K$ and over the closed residue field of $A$.
Over the Artinian coefficient-field algebra $B$, the iterated
argument of D.11 gives a homomorphism
$u_B:\mathbb G_{a,B}\to G_B$ with $u_B(1)=g_B$.
More explicitly, restriction of scalars to $K$ is representable
by a separated finite-type group because $B/K$ is finite free
and a universal homeomorphism. The reduction kernel has a finite
filtration by subgroups of additive vector groups. Its inverse
image over the actual unipotent cyclic image is unipotent;
exponential and adjunction give $u_B$. This uses neither
flatness of $G$ nor surjectivity onto the entire reduced group.

Choose an affine open $V\subset G$ containing the identity at
the closed point. The whole identity section lies in $V$, since
$\Spec A$ is local. The closed-fibre parameter through $g$
meets $V$ on a nonempty open of the affine line. Therefore
$g^n\in V(A)$ for all but finitely many integers $n$: its
closed-point image lies in $V$, and an open of a local spectrum
containing the closed point is the whole spectrum.

Choose finite algebra generators $h_1,\ldots,h_q$ of
$\mathcal O(V)$. Their pullbacks under $u_B$, initially on
$u_B^{-1}(V_B)$, are rational functions $F_i\in B(T)$.
At all but finitely many integers they satisfy
\[
                        F_i(n)=h_i(g^n)\in A.
\]
Apply the primary-ring denominator lemma and take a common
denominator $D\in A[T]$ with nonzero reduction at the closed
point. The coordinate formulas give a map
$v:D(D)\to V$. All defining relations hold because
$A[T,1/D]\hookrightarrow B(T)$ is injective: $A\hookrightarrow B$
is injective and $D$ has nonzero reduction in the residue field
of $B$. The map agrees with $u_B$ on their common domain.

At every prime of the local ring $A$, the polynomial $D$ has
nonzero reduction: a coefficient that is a unit at the maximal
ideal is a unit everywhere. All residue characteristics are
zero. Hence the translates $D(D)+n$, $n\in\mathbb Z$, cover
$\mathbb A^1_A$, because a nonzero polynomial has only finitely
many roots whereas integer translates of any geometric point
are distinct. On $D(D)+n$ put
\[
                         v_n(T)=v(T-n)g^n.
\]
They agree after localization to $B$. That localization is
schematically dense on $\mathbb A^1_A$ and on every open
intersection, since it inverts nonzerodivisors. Separatedness
therefore makes the maps agree, and they glue. The same
schematic-density argument on $\mathbb A^2_A$ proves the
group law, and on $\Spec A$ proves the value at $1$.
This settles every primary local ring.

For a general Noetherian local $\Q$-algebra $A$, take a finite
primary decomposition $0=\bigcap_{i=1}^r Q_i$. Recall briefly
why it exists: the ascending chain condition makes every ideal
a finite intersection of irreducible ideals. An irreducible zero
ideal is primary, since stabilization of the annihilators of
$y^j$ shows that, if $xy=0$ and $x\ne0$, then
$(x)\cap(y^j)=0$ for a sufficiently large $j$, forcing $y^j=0$.
Apply this in the quotients by the irreducible ideals.

The primary-ring argument supplies parameters on all
$A/Q_i$. Choose a common power and rescale. Their restrictions
to every closed overlap agree by D.12, since these are
$\Q$-schemes and the value at $1$ is the same.
Closed patching gives a map over $A$: for two ideals with
intersection zero,
\[
                 A=A/I\times_{A/(I+J)}A/J,
\]
and the same identity holds after adjoining one or two
polynomial variables. Compatible maps from the two closed
pieces to any scheme glue. Indeed, inverse images of an affine
target open form one open in the source, since a finite closed
cover detects openness; the fibre product of structure sheaves
then glues their algebra maps. It also proves uniqueness.
Induct on the number of primary pieces. Patching on the affine
plane checks the group law, and patching the sections checks
the value at $1$. In the induction, the overlap of the next
primary piece with the already patched union is itself a
$\Q$-scheme; D.12 makes the two parameters equal on this full
scheme-theoretic overlap. Pairwise agreement on reduced
intersections alone is not used as a descent criterion.

Finally, over a general Noetherian $\Q$-scheme, spread the local
maps, their group laws, and their values from local rings to
open neighbourhoods by finite presentation. Take a finite
cover and a common multiple of its exponents. Characteristic-zero
uniqueness in D.12 makes the rescaled maps agree on overlaps,
so they glue. That same result supplies the asserted uniqueness.
\end{proof}

\begin{statement}{D.19}{Proved}{Integer localization with torsion errors}
Let $A$ be a local ring, $G/A$ separated of finite
presentation, and $g\in G(A)$ fibrewise quasi-unipotent.
If some positive power of $g$ has an additive parameter over
$A[1/n]$, for an integer $n>0$, then some positive power has
an additive parameter over $A$. The integer $n$ may be a
zerodivisor in $A$.
\end{statement}
\begin{proof}
If $n$ is a unit there is nothing to prove; if it is nilpotent,
D.13 applies. Otherwise let $u$ over $A[1/n]$ have value
$g^m$ at $1$. Choose an affine identity chart $V=\Spec C$
as in D.16. This affine open is of finite presentation over $A$:
its inclusion in the quasi-separated finite-presentation scheme
$G$ is a quasi-compact open immersion. The inverse image of $V$ contains the whole zero
section; choose $f$ with $f(0)=1$ and $D(f)$ in that inverse
image. Write the finitely many coordinate images as
$a_i/f^{k_i}$, with constant terms $e^*(h_i)\in A$.
After replacing $T$ by $n^N T$, choose lifts
$F,A_i\in A[T]$ with
\[
          F(0)=1,\quad F\equiv1\pmod n,\quad
          A_i(0)=e^*(h_i).
\]
The algebra $C$ is finitely presented. Evaluate its finite set
of defining relations at $A_i/F^{k_i}$, and clear powers of
$F$. Their numerators $P_j(T)$ vanish after inverting $n$
and have constant term zero. Each is killed by a power of $n$.
Here $A[T]\hookrightarrow A[T,1/F]$ is injective: a polynomial
with constant term $1$ is a nonzerodivisor, as comparison of
the lowest nonzero coefficient shows. A further substitution
$T\mapsto n^eT$ therefore kills all the relation numerators.

We obtain a morphism $D(F)\to V$, agreeing with a rescaling of
$u$ after inverting $n$. The cover
$D(F)\cup\mathbb A^1_{A[1/n]}=\mathbb A^1_A$ glues it
to a scheme map $\phi:\mathbb A^1_A\to G$ with
$\phi(0)=e$ exactly. It need not yet be a homomorphism.

The equalizer of $\phi(T+U)$ and $\phi(T)\phi(U)$ is a
finitely presented closed subscheme of $\mathbb A^2_A$.
Its ideal is finitely generated, becomes zero after inverting
$n$, and every generator has zero constant coefficient.
Likewise for the equalizer of $\phi(T)g^m$ and $g^m\phi(T)$:
after inverting $n$, a homomorphism commutes with its value
$g^m$ at $1$. A common further dilation by a power of $n$
kills all the coefficient errors. The resulting map $\psi$
is a group homomorphism and centralizes $g^m$. Put
\[
                  h=\psi(1),\qquad
                  a=mn^b,\qquad\delta=h^{-1}g^a,
\]
where $b$ is the total dilation exponent. Then $h$ and $g^a$
commute, and $\delta=1$ after inverting $n$.

Every geometric fibre of $\delta$ has finite order. Outside
$V(n)$ it is the identity. On $V(n)$ the residue characteristic
is some $p\mid n$; there $h$ is killed by $p$ and $g^a$ has
finite order by quasi-unipotence. Their commuting product has
finite order. Constructible compactness, applied to the
finitely presented equalizers of $\delta^q$ and $1$, gives
$q>0$ such that $\delta^q$ is the identity in every geometric
fibre, exactly as in D.13.

Let $I$ be its equalizer ideal. It is finitely generated and
nilpotent, say $I^r=0$. It also vanishes after inverting $n$,
so $n^c I=0$ for some $c\ge0$. The successive square-zero
identity kernels for $I^j/I^{j+1}$ are additive and killed
by $n^c$. Consequently
\[
                         \delta^{q n^{c(r-1)}}=1.
\]
Write this exponent as $d$. Because $h$ and $g^a$ commute,
$g^{ad}=h^d=\psi(d)$. Thus $T\mapsto\psi(dT)$ is the required
parameter. This handles the value error as well as the relation
and group-law errors; generic agreement alone would not do so
over a ring with integer torsion.
Every ideal used in this proof has an explicit finite set of
generators, supplied by finite presentation. A finitely generated
nil ideal is nilpotent over any ring. Thus no Noetherian hypothesis
on $A$ is needed in this localization lemma.
\end{proof}

\begin{statement}{D.20}{Proved}{The positive-power theorem over every Noetherian base}
Let $S$ be a Noetherian scheme, let $G/S$ be a separated group
scheme of finite presentation, and let $g\in G(S)$ be fibrewise
quasi-unipotent in the sense of \S50.1, C4 scan 5. Then there
exist an integer $m>0$ and a homomorphism
\[
                 u:\mathbb G_{a,S}\longrightarrow G,
                 \qquad u(1)=g^m.
\]
No reducedness, irreducibility, regularity, or flatness assumption
on $S$ is required, and $G$ need not be affine, smooth, or flat.
\end{statement}
\begin{proof}
Apply D.18 over $S_{\Q}$. Finite-presentation limit descent
of the map, group law, and value at $1$ supplies a parameter
for some power over $S[1/n]$, for one integer $n>0$, as in
D.16. The inverse system has affine transition maps and
Noetherian, hence quasi-compact quasi-separated, bases.

At every local ring of $S$, D.19 now supplies a positive-power
parameter. Spread it and its identities to an open neighbourhood.
A finite cover exists, and the powers can be made common.
D.12 descends the local maps after one further common dilation.
This proves the assertion. The argument includes the case
$S_{\Q}$ is empty.
\end{proof}

\begin{statement}{D.21}{Proved}{All quasi-compact quasi-separated bases}
Let $S$ be a quasi-compact quasi-separated scheme, $G/S$ a
separated group scheme of finite presentation, and $g\in G(S)$
fibrewise quasi-unipotent. There are $m>0$ and a homomorphism
$u:\mathbb G_{a,S}\to G$ with $u(1)=g^m$.
Thus the positive-power assertion of \S50.2, C4 scans 10--13,
is true with this base hypothesis. The base may be non-Noetherian;
the group need not be affine, smooth, or flat.
\end{statement}
\begin{proof}
We first treat an arbitrary $\Q$-algebra $R$. Descend $G$, its
group operations, and $g$ to a finitely generated $\Q$-subalgebra
$R_0\subset R$. After enlarging $R_0$ the descended group $G_0$
is separated and of finite presentation. These assertions use
\citet[Lemma 32.10.1, Tag 01ZM]{DN_StacksLimits} for objects,
morphisms, and identities, and
\citet[Lemma 32.8.6, Tag 01ZQ]{DS_StacksSeparatedLimit} for
separatedness.

Let $C\subset\Spec R_0$ be the geometric fibrewise
quasi-unipotence locus of $g_0$. It is a union of constructible
subsets for the following elementary reason. Given $x\in C$,
a homomorphism over $\overline{\kappa(x)}$ taking $1$ to
$g_0^m$ descends, together with its group law and value, to a
finitely generated $R_0$-subalgebra
$B\subset\overline{\kappa(x)}$. Thus $\Spec B\to\Spec R_0$
has constructible image by Chevalley's theorem
\citep[Theorem 29.23.3, Tag 054K]{DS_StacksChevalley}.
Every point in this image belongs to $C$: after a residue-field
extension $L/k$ the same homomorphism is a witness. Its two
pullbacks over $L\otimes_k L$ agree by D.12, since this is
a $\Q$-algebra and both have the same value at $1$. Faithfully
flat descent of morphisms therefore gives a parameter over $k$
itself. Conversely every $x\in C$ is in one
of these constructible images. This proves the assertion about
$C$ without representing the full functor of homomorphisms.

Their inverse images cover $\Spec R$, by the hypothesis on $g$.
They are open in its constructible topology, so compactness
\citep[Lemma 5.23.2, Tag 0901]{DN_StacksConstructible} gives
a finite subcover. Write $C'$ for the corresponding finite
union of constructible subsets of $\Spec R_0$. The complement
$E=\Spec R_0\setminus C'$ has empty inverse image in $\Spec R$.
Since $R_0$ is Noetherian, write $E$ as a finite union of sets
$V(I_\alpha)\cap D(f_\alpha)$ with each $I_\alpha$ finitely
generated. Empty inverse image means
\[
                      (R/I_\alpha R)[1/f_\alpha]=0.
\]
The finitely many identities expressing this vanishing use
only finitely many coefficients from $R$. A finitely generated
$\Q$-subalgebra $R_1\subset R$ containing these coefficients
and $R_0$ therefore has its entire spectrum mapped into $C'$.
Every geometric fibre of $(G_0,g_0)_{R_1}$ is quasi-unipotent.
Apply D.18 over the Noetherian ring $R_1$, and pull the resulting
parameter back to $R$. This proves the assertion for every
affine $\Q$-base. A finite affine cover and D.12 give it over
every quasi-compact quasi-separated $\Q$-scheme.

For general $S$, apply that result to $S_{\Q}$. The inverse
system $S[1/n]$ has affine transition maps and quasi-compact
quasi-separated terms. Finite-presentation limit descent
supplies a parameter through a power of $g$ over some $S[1/n]$.
D.19 applies at each local ring of $S$, with no Noetherian
assumption. Spreading each resulting parameter and its identities
to a neighbourhood, choosing a finite cover, and applying D.12
after a common dilation proves the assertion.
\end{proof}

\begin{statement}{D.22}{Known}{Smooth integral models of compact adelic representations}
Let $H/\Q$ be an affine group scheme of finite type, $M$ a
profinite group, $\rho:M\to H(\mathbb A_f)$ a continuous
homomorphism, and $\chi:H\to\mathbb G_{m,\Q}$ a character.
There exist a smooth affine group scheme $\mathscr H/\mathbb Z$
of finite type, an identification $\mathscr H_{\Q}=H$, a
character $\chi_{\mathbb Z}:\mathscr H\to\mathbb G_{m,\mathbb Z}$
extending $\chi$, and a continuous factorization
\[
                  M\longrightarrow\mathscr H(\widehat{\mathbb Z})
                    \longrightarrow H(\mathbb A_f)
\]
of $\rho$. If a subgroup $\Lambda\subset M$ has its image
in diagonal $H(\Q)$, then that image lies in $\mathscr H(\mathbb Z)$.

This answers the smooth-model existence question of \S52.3,
folio 754, C4 scan 213, with the adelic factorization specified
in (45) and (47), scans 211--212. It does not assert the later
supplementary conditions (55)--(64). No connectedness hypothesis
on $H$ is required.
\end{statement}
\begin{proof}
Choose a faithful closed representation $H\hookrightarrow
\operatorname{GL}(V)$ over $\Q$
\citep[Theorem 4.9 and Corollary 4.10]{D_Milne2022}.
We explain the integral lattice
step. The compact subgroup $K=\rho(M)$ maps into
$\operatorname{GL}(V)(\mathbb A_f)$. For every prime $p$, its
projection $K_p$ is compact, hence bounded. If $L_{0,p}$ is a
lattice in $V_{\Q_p}$, the $\mathbb Z_p$-module
\[
                  L_p=\sum_{k\in K_p}kL_{0,p}
\]
lies between $L_{0,p}$ and $p^{-a_p}L_{0,p}$ for some $a_p$.
It is therefore a lattice, and it is stable under $K_p$.
Compactness in the restricted product implies that
$K_p\subset\operatorname{GL}(L_{0,p})$ for all but finitely
many $p$, with $L_{0,p}$ taken from one rational lattice $L_0$.
At those primes take $L_p=L_{0,p}$.

These local lattices come from a global lattice
$L\subset V$. For completeness, choose $N>0$ with
$NL_{0,p}\subset L_p\subset N^{-1}L_{0,p}$ for all $p$.
The finite local submodules of
$N^{-1}L_{0,p}/NL_{0,p}$ specify, by the Chinese remainder
theorem, a subgroup of $N^{-1}L_0/NL_0$; its inverse image is
$L$, and $L\otimes\mathbb Z_p=L_p$.
The character of a compact subgroup has valuation zero at every
prime: the image of the continuous homomorphism
$v_p\circ\chi:K_p\to\mathbb Z$ is finite, hence zero.
Thus $\chi(K)\subset\widehat{\mathbb Z}^{\times}$.

Take the schematic closure $\mathscr H_0$ of the graph of
$\chi$ in $\operatorname{GL}(L)\times\mathbb G_{m,\mathbb Z}$.
It is an affine group scheme of finite type, flat over
$\mathbb Z$, with generic fibre $H$; flatness follows because
its coordinate ring is torsion-free. Projection to the second
factor extends $\chi$. The whole subgroup $K$ lies in
$\mathscr H_0(\widehat{\mathbb Z})$: at every prime its
matrix and character coordinates are integral, and the closure
equations vanish because they vanish on the generic fibre.

Cartier's theorem makes $H$ smooth
\citep[Lemma 39.8.2, Tag 047N]{D_StacksCartier}, and finite presentation
shows that $\mathscr H_0$ is smooth away from finitely many
primes. At each remaining prime apply the group-smoothening
theorem of \citet[\S7.1, Theorem 5, p.~175]{DS_BLR1990}
over $R=\mathbb Z_{(p)}$. It supplies a smooth finite-type
group $\mathscr H_p'\to(\mathscr H_0)_R$, isomorphic on
generic fibres, through which every morphism from a smooth
$R$-scheme factors uniquely. Its construction is a finite
sequence of affine dilatations, so $\mathscr H_p'$ is affine
\citep[\S3.2, pp.~62--64; \S7.1, pp.~174--175]{DS_BLR1990}.

We check explicitly that all $\mathbb Z_p$-points are retained.
The extension $\mathbb Z_{(p)}\subset\mathbb Z_p$ has
ramification index one and the same residue field.
Given $a\in\mathscr H_0(\mathbb Z_p)$,
\citet[\S3.6, Proposition 4, p.~79]{DS_BLR1990} supplies
a finite sequence of dilatations $X\to(\mathscr H_0)_R$
such that $a$ lifts into the $R$-smooth locus of $X$.
That locus maps to $\mathscr H_p'$ by the preceding universal
property, yielding a lift of $a$. The lift is unique by the
generic-fibre isomorphism and separatedness. This argument
does not assume that the original group is a N\'eron model.

Spread these finitely many local modifications to neighbourhoods
of their primes and glue them with the unchanged smooth model
on the complementary open. The modifications are generic
isomorphisms, so after shrinking each neighbourhood they agree
on its punctured part. This gives a smooth affine finite-type
$\mathbb Z$-group $\mathscr H\to\mathscr H_0$ with unchanged
generic fibre. Affineness is local on the base. This is also
the finite-local-modification gluing used in
\citet[proof of \S7.1, Corollary 6, p.~175]{DS_BLR1990}.
Compose with $\mathscr H_0\to\mathbb G_m$ to extend $\chi$.
The equality on all $\mathbb Z_p$-points supplies the required
factorization. It is continuous: for an affine finite-type
$\mathbb Z_p$-scheme its set of $\mathbb Z_p$-points is compact,
so the continuous bijection onto the original point set is a
homeomorphism; take the product over $p$.

Finally, for any affine $\mathbb Z$-scheme with the specified
generic fibre,
\[
             H(\Q)\cap\mathscr H(\widehat{\mathbb Z})
                         =\mathscr H(\mathbb Z).
\]
Indeed each coordinate of a rational point integral at every
prime belongs to $\Q\cap\prod_p\mathbb Z_p=\mathbb Z$.
This proves the assertion for $\Lambda$ and gives (47) for
the prescribed rational images of $\Pi_0$ and $\mathcal G_0$.
Their schematic closures contain these integral sections. If
one of the quotient schemes in (46) is represented, the
character, trivial on the relevant closure by generic density,
factors through it as in (49). Representability remains the
explicit proviso in the manuscript, not an extra conclusion
silently used here.
\end{proof}

\begin{statement}{D.23}{Refuted}{The integral-point recap and uniqueness from the preceding data}
The literal replacement of $M\to\mathscr H(\widehat{\mathbb Z})$
by $M\to\mathscr H(\mathbb Z)$ in the parenthetical recap
on C4 scan 213 is false for a Lie homomorphism whose multiplier
has open image in $\widehat{\mathbb Z}^{\times}$, as required
in \S52.3, scans 201 and 207. The factorization in (45) and
(47) must be kept distinct from this stronger integral-point
reading.

Moreover, the existence data (45)--(49) alone do not specify
a unique smooth model with its fixed generic-fibre identification.
The author's suggestion of uniqueness under further conditions
does not state such conditions and is not treated as a uniqueness
theorem.
\end{statement}
\begin{proof}
If the whole of $M$ mapped to $\mathscr H(\mathbb Z)$ and
$\chi$ extended over $\mathbb Z$, its multiplier would take
values in $\mathbb G_m(\mathbb Z)=\{\pm1\}$. This contradicts
its open image in $\widehat{\mathbb Z}^{\times}$, since no
finite subgroup is open there. For the last assertion, every
basic neighbourhood imposes congruence conditions at only
finitely many primes and leaves an infinite unit subgroup at
another prime. Already $H=\mathbb G_m$, $M=\widehat{\mathbb Z}^{\times}$,
$\rho$ the natural inclusion, and $\chi=\mathrm{id}$ give
the explicit obstruction.

To see the failure of uniqueness from the preceding data, use
\[
 \begin{aligned}
 H&=\mathbb G_a\rtimes\mathbb G_m,
 & (x,u)(y,v)&=(x+uy,uv),\\
 M&=\widehat{\mathbb Z}\rtimes\widehat{\mathbb Z}^{\times}.
 \end{aligned}
\]
with its natural adelic representation and character $(x,u)\mapsto u$.
For each integer $c>0$, take the smooth model
$\mathscr H_c=\mathbb G_{a,\mathbb Z}\rtimes\mathbb G_{m,\mathbb Z}$
whose generic identification with $H$ sends $(z,u)$ to $(z/c,u)$.
The adelic factorization is $(a,u)\mapsto(ca,u)$, and the
character extends. Rational integral subgroups, such as the
integer translations, satisfy the required factorization too.
An isomorphism $\mathscr H_1\simeq\mathscr H_2$ inducing the
identity on $H$ would have to be $(z,u)\mapsto(2z,u)$ by generic
density. This morphism is not an isomorphism at $2$, where it
kills the additive coordinate. Thus these are distinct models
of the specified rational group, although the underlying
unmarked group schemes are isomorphic. This does not refute
a uniqueness theorem with additional precisely stated conditions;
it proves that the data currently preceding scan 213 are
insufficient for one.
\end{proof}

\paragraph{Scope after the surface argument.}
D.21 settles every quasi-compact quasi-separated base for separated
finitely presented groups. D.14 proves that quasi-compactness
alone, the literal base hypothesis in the manuscript, is
insufficient. Neither Noetherian singularities and nilpotents
nor non-Noetherian affine bases are retained as unresolved tests.
The smooth integral-model existence question at scan 213, folio 754,
with the coherent adelic reading of (45)--(49), is answered by D.22;
the stronger integral-point recap and bare uniqueness are corrected
in D.23. Conditions (55)--(64), introduced later
as additional hypotheses and then used to define a nondegenerate
integral approximation, are not an asserted universal existence
theorem.

The argument does not infer representability of the full functor
$\operatorname{Hom}_{\rm gp}(\mathbb G_a,G)$. Such a shortcut
would already fail for $G=\mathbb G_m$ over a characteristic-zero
field: nilpotent test rings admit nontrivial exponential maps,
and this Hom functor is not represented by a scheme
\citep[p.~2]{DS_BrionHom}. The proof uses coordinate denominators,
integer interpolation, and finite nilpotent filtrations instead.
